\subsection{Comparison with existing methods}

\paragraph{Quality criteria and per-point tools.} Rank-based quality criteria
\citep{venna2001,kaski2003,chen2009,lee2009,venna2010} and per-point tools such as scDEED
\citep{xia2024} and EMBEDR \citep{johnson2022} measure how well a map agrees with the data it was
made from. With trial-averaged data they measure agreement with the averages, and they inherit the
averages' noise. scDEED's permutation null, and EMBEDR's embeddings of randomized data, answer
whether a neighborhood is better preserved than chance, which is a different question from
whether it is faithful to the noise-free responses. In both datasets, per-point trustworthiness and
scDEED's trustworthy label left most false comfort, and scDEED passed more EEG neighborhoods than
the plug-in check without being more often right (Table~\ref{tab:checks}).

\paragraph{Stability.} Recurrence across t-SNE runs, bootstrap resamples (as in DynamicViz;
\citealp{sun2023}) and split halves measures whether a neighborhood is stable, not whether it is
faithful. A neighborhood created by the noise in the averages recurs in every map of those
averages, and a map's systematic distortions recur across seeds and resamples. Seed recurrence
passed a third of the EEG neighborhoods, most of them not faithful. The strict resampling checks
behaved differently: at 8 trials they passed almost nothing, which is itself a useful warning, and
at 32 trials in EEG what they passed was more often faithful. Among the checks we tested, only a
strict resampling check avoided passing many unfaithful neighborhoods, but at 8 trials even
split-half recurrence was wrong about roughly half of the little it passed.

\paragraph{Cross-validated distances and noise ceilings.} RSA's cross-validated distances remove
the positive bias of plug-in distances \citep{walther2016,guggenmos2018}, and noise ceilings bound
how well any model can match a dissimilarity matrix \citep{nili2014,schutt2023}. Neither is
designed to certify single neighborhoods. Cross-validated distances made faithful calls rare, both
as the map's distances (Section~\ref{sec:res-rsa}) and as a check on a t-SNE map, but the calls that
remained were still mostly wrong.

\paragraph{Why checks computed from the same trials cannot see the error.} Under Gaussian trial
noise, the averages carry all the information that the trials hold about the noise-free means; the
residuals add information only about the size of the noise (Supplementary Section~S1.3). Split
halves are the averages plus and minus a noise-only difference, bootstrap resamples are the averages
plus resampled noise, and cross-validated squared distances are the plug-in ones minus a noise-only
term. Every check of Section~\ref{sec:checks} is therefore a function of the averages, the noise
level and randomization. It can be stricter than the plug-in check, but it cannot tell where the
averages themselves are wrong. We recorded this argument, and the prediction that all six checks
would leave most false comfort, in the comparison's protocol before it was run. The prediction held
for the four checks that passed many neighborhoods. The two strict checks did better only by
passing almost nothing: strictness selects the neighborhoods whose margins are large relative to
the noise, and at 8 trials almost none exist. Per-point certificates with family-wise error
control are possible in principle, for instance through simultaneous inference on ranks
\citep{mogstad2024}. In our pilot on THINGS-EEG2, certificates of this kind, built from simultaneous
studentized bounds on pairwise distances, resolved no neighborhood even at 72 trials
(Supplementary Section~S1.5), consistent with Table~\ref{tab:snr}.

\paragraph{Denoised and deeper references.} \citet{lin2024} also judge maps against a denoised
signal, obtained by PCA truncation of the same data, to choose hyperparameters for a whole map. Here
the reference comes from held-out repetitions, it is used at the level of individual neighborhoods,
and false comfort is measured as a function of the number of trials. Downsampling studies compare
neighbor graphs from shallow data with graphs from deeper data \citep{ahlmann2023}, treating the
deeper data as the truth; our extrapolation removes the reference's own noise instead.

\paragraph{What the estimator adds.} The estimator needs held-out repetitions and gives a
dataset-level answer: how much of what a check passes is wrong, with a conservative bound. It does
not certify individual neighborhoods, which the certificates of our pilot could not do at these
signal-to-noise ratios either. It tells an analyst whether reading a map's neighborhoods is
warranted at all, and how false comfort falls as trials are added.
